\chapter{Examples of coarea, multiplicity and projection}\label{sec:examples-coarea-projections}
These calculations illustrate the area and coarea conventions in \cite{Sim18}, Chapter 2, Sections 3 and 7, and Chapter 3, Sections 2--3. Each example specifies whether it counts points, measures a positive-dimensional fiber, or pushes an unweighted surface measure forward. Exceptional levels can carry large fibers even when every almost-everywhere conclusion is valid.

\begin{example}[A tilted plane checks the coarea metric factor]\label{ex:coarea-tilted-plane}
Let $M_a=\{(s,t,as):(s,t)\in(0,1)^2\}\subset\mathbb R^3$ and $f(x,y,z)=x$, with $a\in\mathbb R$. Then
\[
 d\mathcal H^2\llcorner M_a=\sqrt{1+a^2}\,ds\,dt,
 \qquad J_1(d^{M_a}f)=\frac1{\sqrt{1+a^2}}.
\]
For $0<u<1$, the fiber $M_a\cap f^{-1}(u)$ has length one. Thus coarea gives $\int_{M_a}J_1(d^{M_a}f)\,d\mathcal H^2=1$, whereas the unweighted area is $\sqrt{1+a^2}$.
\end{example}
\begin{proof}
The graph derivative has tangent columns $(1,0,a)$ and $(0,1,0)$, which are orthogonal with squared lengths $1+a^2$ and one. Its area Jacobian is therefore $\sqrt{1+a^2}$. In the corresponding orthonormal tangent basis, the derivative of $f$ is the row $(1/\sqrt{1+a^2},0)$, giving the stated coarea Jacobian. A fiber has fixed $s=u$ and varies only in $t$, with unit speed and length one. Integrating either side of coarea gives one. Replacing the tangential Jacobian by the ambient gradient norm $|Df|=1$ would incorrectly give the larger unweighted area when $a\ne0$.
\end{proof}

\begin{example}[Height on a round two-sphere has a constant area density]\label{ex:coarea-sphere-height}
Let $R>0$ and let $z$ denote height on the sphere $S_R^2\subset\mathbb R^3$. Away from its poles,
\[
 |\nabla^{S_R^2}z|=\frac{\sqrt{R^2-z^2}}R,
 \qquad \mathcal H^1(S_R^2\cap\{z=t\})=2\pi\sqrt{R^2-t^2}
 \quad(-R<t<R).
\]
The pushforward of area by height is $2\pi R\,\chi_{(-R,R)}\mathcal L^1$. In particular the sphere has area $4\pi R^2$, and
\[
 \int_{S_R^2}|\nabla^{S_R^2}z|\,dA=\pi^2R^2.
\]
\end{example}
\begin{proof}
The unit normal is $x/R$. Projection of $e_3$ onto the tangent plane has norm $\sqrt{1-(z/R)^2}$, proving the gradient formula. The fiber is a horizontal circle of the indicated radius. The inverse-Jacobian formula in Proposition \ref{prop:coarea-inverse-jacobian} gives the area pushforward density as the fiber length divided by this gradient norm, namely $2\pi R$. The two poles have zero area and contribute no critical pushforward mass. Integrating the constant density over $(-R,R)$ gives the area. Ordinary coarea gives the last integral as
\[
 2\pi\int_{-R}^R\sqrt{R^2-t^2}\,dt=\pi^2R^2,
\]
since the integral is the area of a semicircle. Equivalently the parametrization $(\sqrt{R^2-t^2}\cos\theta,\sqrt{R^2-t^2}\sin\theta,t)$ has orthogonal derivatives with lengths $\sqrt{R^2-t^2}$ and $R/\sqrt{R^2-t^2}$, and hence constant area Jacobian $R$. The divergent reciprocal near the poles cancels the shrinking fiber length.
\end{proof}

\begin{example}[Squared radius on a disk gives a uniform pushforward]\label{ex:coarea-squared-radius-disk}
Let $R>0$, $D_R=\{x\in\mathbb R^2:|x|<R\}$ and $f(x)=|x|^2$. Then
\[
 f_\#(\mathcal L^2\llcorner D_R)=\pi\chi_{(0,R^2)}\mathcal L^1.
\]
For $0<t<R^2$, the fiber length is $2\pi\sqrt t$ and $|Df|=2\sqrt t$ on the fiber. The origin is critical but has zero area.
\end{example}
\begin{proof}
The derivative is $2x$, so the only critical point is the origin. Each positive fiber in the range is a circle of radius $\sqrt t$. Proposition \ref{prop:coarea-inverse-jacobian} gives density $(2\pi\sqrt t)/(2\sqrt t)=\pi$ for the unweighted pushforward. There is no atom from the critical point because it has zero two-dimensional measure. Independently, the sublevel set $\{f<t\}$ has area $\pi t$ for $0<t<R^2$, whose derivative is the same constant. Coarea with its Jacobian weight instead gives density $2\pi\sqrt t$; the two pushforwards must not be conflated.
\end{proof}

\begin{example}[A flat plateau produces a critical atom]\label{ex:coarea-flat-plateau-atom}
On $D=(-1,1)\times(0,1)\subset\mathbb R^2$, let $f(x,y)=\max\{x,0\}$. Its unweighted area pushforward is
\[
 f_\#(\mathcal L^2\llcorner D)=\delta_0+\chi_{(0,1)}\mathcal L^1.
\]
Its Jacobian-weighted pushforward is only $\chi_{(0,1)}\mathcal L^1$. The exceptional fiber $f^{-1}(0)$ has infinite $\mathcal H^1$ measure, although almost every fiber in the range has length one.
\end{example}
\begin{proof}
On $x<0$ the map is constant and its gradient is zero. This half-rectangle has area one and pushes to a unit atom at zero. On $x>0$ the gradient is $(1,0)$, and Fubini gives Lebesgue measure on $(0,1)$ as the pushforward. The line $x=0$ is a null set for area, so its nondifferentiability has no effect. Multiplication by the Jacobian removes the flat half-rectangle and yields the second pushforward.

The fiber at zero contains a nonempty open rectangle. Its one-dimensional Hausdorff measure is infinite: if a bounded set had finite $\mathcal H^1$, diameter covers would give zero two-dimensional Lebesgue measure by multiplying their one-dimensional cost by the maximum diameter and letting that diameter tend to zero. This contradicts the positive area of the half-rectangle. For $t\in(0,1)$ the fiber is a vertical unit segment, and other nonzero levels have empty fibers. An infinite fiber value at one Lebesgue-null level does not affect the coarea integral. This example shows why the critical contribution in Proposition \ref{prop:coarea-inverse-jacobian} must be retained.
\end{proof}

\begin{example}[Distance is constant on a sphere and its volume function jumps]\label{ex:coarea-spherical-critical-level}
Let $R>0$, $M=S_R^2\subset\mathbb R^3$ and $\rho(x)=|x|$. Its tangential gradient vanishes everywhere, while
\[
 \rho_\#(\mathcal H^2\llcorner M)=4\pi R^2\delta_R.
\]
The function $V(t)=\mathcal H^2(M\cap B_t(0))$ equals zero for $t\leq R$ and $4\pi R^2$ for $t>R$. It is not absolutely continuous. Coarea nevertheless holds: all level intersections except the single level $R$ are empty.
\end{example}
\begin{proof}
The function $\rho|M$ is the constant $R$, so its tangential derivative is zero and its unweighted pushforward is the claimed atom. The area is $4\pi R^2$ by Example \ref{ex:coarea-sphere-height}. The strict inequality in the open-ball definition accounts for $V(R)=0$. At the one nonempty level the fiber is the entire two-sphere and has infinite one-dimensional measure, by the same dimension comparison as in Example \ref{ex:coarea-flat-plateau-atom}, applied to a regular chart. That value occurs on a null set of radii, and the integral of the fiber length is zero, agreeing with the zero Jacobian integral. Proposition \ref{prop:rectifiable-distance-volume-derivative} concerns the empty regular part in this example and makes no incorrect absolute-continuity claim for $V$ itself.
\end{proof}

\begin{example}[A fold counts two preimages rather than one image]\label{ex:area-fold-multiplicity}
For $f:(-1,1)\to\mathbb R$, $f(s)=s^2$, the image has length one, but
\[
 \int_{-1}^1|f'(s)|\,ds=2
 =\int_0^1\#f^{-1}(t)\,dt.
\]
The unweighted pushforward of $ds$ has density $t^{-1/2}\chi_{(0,1)}(t)$, which is integrable despite its singularity at zero. This is an equal-dimensional area-formula example, not a positive-dimensional coarea fiber.
\end{example}
\begin{proof}
For $0<t<1$ there are two preimages $\pm\sqrt t$, and outside $[0,1)$ there are none. The unique preimage at zero does not affect the integral. Direct integration gives $\int_{-1}^1 2|s|\,ds=2$, matching the multiplicity integral. For the unweighted pushforward, apply the weighted area formula with weight $1/|f'|$ away from the null point zero. At each positive target value the two inverse derivative factors are each $1/(2\sqrt t)$, summing to $1/\sqrt t$. Its integral over $(0,1)$ is two, the original interval length. Thus counting the image once would lose one sheet, while omitting the inverse derivative would describe a different pushforward measure.
\end{proof}

\begin{example}[A segment and a separated dust realize both decomposition parts]\label{ex:rectifiable-part-segment-dust}
Let $L=[0,1]\times\{0\}$ and let $K'=K+(0,2)$ be a translated four-corner set. For $A=L\cup K'$, its rectifiable part is $L$ and its purely one-unrectifiable part is $K'$, modulo $\mathcal H^1$-null sets. Both parts have positive finite length.
\end{example}
\begin{proof}
The segment is a compact Lipschitz image of an interval and has length one. The translated four-corner set has positive finite length and is purely one-unrectifiable by Theorem \ref{thm:four-corner-pure-unrectifiability}. Translation preserves those properties. The two sets are disjoint by their height ranges, so they already give a decomposition of the required type. Uniqueness in Theorem \ref{thm:rectifiable-pure-decomposition} identifies them with the canonical measure classes of the two parts. In particular no positive-length portion of the dust can be absorbed into a different rectifiable part by changing the chosen representative.
\end{proof}

\begin{example}[One null coordinate projection does not imply pure unrectifiability]\label{ex:one-null-projection-insufficient}
For $1\leq m<n$, consider
\[
 M=\operatorname{span}\{e_1,\ldots,e_{m-1},e_{m+1}\}\subset\mathbb R^n,
\]
with the empty initial list allowed when $m=1$. Projection onto $\operatorname{span}\{e_1,\ldots,e_m\}$ has $m$-dimensional image measure zero, although $M$ is a smooth rectifiable $m$-plane. Projection onto its own coordinate plane is the identity.
\end{example}
\begin{proof}
The first projection maps $M$ to the $(m-1)$-plane spanned by $e_1,\ldots,e_{m-1}$. This lower-dimensional plane has zero $\mathcal H^m$ measure, as follows from diameter covers of bounded pieces. Yet the plane $M$ itself has positive $m$-measure on every nonempty relative ball, and is countably rectifiable by its linear parametrization. It is therefore not purely $m$-unrectifiable. The second coordinate projection has Jacobian one on $M$. This explains why Theorem \ref{thm:null-coordinate-projections-pure} requires all coordinate $m$-projections, including when $n<2m$ so a fully orthogonal $m$-plane would not exist.
\end{proof}
